\chapter{First Aid Kits}

\section*{Definition and Overview}
A first aid kit is a collection of supplies and equipment kept ready for treating injuries and managing emergencies until professional help arrives. A well-stocked kit contains the basics for cuts, grazes, burns, sprains, and minor illnesses, and is organised so that items can be found quickly in an emergency. Kits should be suited to their setting: a home kit differs from a car kit, and workplace kits follow specific standards. A kit is only useful if it is complete, in date, and known to the household, so regular checks and basic first aid knowledge matter as much as the contents. This chapter explains what to put in a first aid kit, how to maintain it, and how to use it.

\section*{Epidemiology}
Most injuries treated with first aid happen in and around the home, where the majority of families keep a first aid kit of some kind. However, many home kits are incomplete, contain expired items, or are not easily found when needed. Workplace first aid kits are mandatory in Australia and must meet regulated standards, while car kits are recommended but less commonly kept. The most common items used are dressings, bandages, and antiseptics for cuts and grazes, followed by treatments for burns, sprains, and headaches. Kit contents reflect the injuries most likely in each setting, such as snake-bite bandages in rural areas.

\section*{Aetiology and Pathophysiology}
The kit contents address the common mechanisms of injury.
\begin{itemize}
    \item \textbf{Cuts and grazes:} dressings, antiseptic wipes, and adhesive bandages clean and protect wounds and reduce infection
    \item \textbf{Bleeding:} sterile gauze and pressure bandages control bleeding by direct pressure
    \item \textbf{Burns:} cool running water is the first treatment; cling film and sterile dressings cover the burn afterwards
    \item \textbf{Sprains and strains:} compression bandages and cold packs limit swelling and support the joint
    \item \textbf{Splinters and foreign bodies:} tweezers and sterile needles remove them cleanly
    \item \textbf{Fever and pain:} paracetamol and ibuprofen manage common symptoms while awaiting care
    \item \textbf{Allergic reactions:} antihistamines and an adrenaline autoinjector (EpiPen) where prescribed
    \item \textbf{Protection for the first aider:} gloves and a resuscitation mask prevent infection transmission
    \item \textbf{Emergency information:} first aid booklets, contact numbers, and medical information forms guide the response
\end{itemize}

\section*{Clinical Features}
A good first aid kit is recognisable by its organisation and completeness.
\begin{itemize}
    \item A sturdy, labelled container that is waterproof and easy to carry
    \item Clearly organised sections or a contents list, so items are found quickly
    \item Items appropriate to the setting: home, car, workplace, or outdoor activities
    \item In-date medicines and sterile items, with expired stock replaced
    \item A first aid guide or instruction card inside
    \item Emergency phone numbers, including 000 and the Poisons Information Centre
    \item A record of the kit's location and any household medical information
    \item For workplaces: compliance with the relevant code of practice
\end{itemize}

\section*{Differential Diagnosis}
Kit contents should match the likely injuries.
\begin{itemize}
    \item Home kit: dressings, bandages, antiseptic, thermometer, pain relief, and basic equipment
    \item Car kit: additional items such as a blanket, torch, emergency triangle, and reflective vest
    \item Workplace kit: regulated minimum contents plus items for the workplace's specific hazards
    \item Sports and outdoor kits: cold packs, sports tape, blister treatment, and snake-bite bandages in Australia
    \item Personal kits: medicines for individuals, including asthma puffers, adrenaline autoinjectors, and insulin
    \item The need to match the kit to the setting rather than buying a single generic kit for everything
\end{itemize}

\section*{When to Seek Medical Care}
A first aid kit treats minor injuries and stabilises serious ones, but it never replaces professional care. Any serious injury, deep wound, burn larger than the palm, head injury, or condition that does not improve needs medical review, and life-threatening situations need 000.

\textbf{Call 000 (or local emergency services) for:}
\begin{itemize}
    \item Severe bleeding, unconsciousness, or breathing difficulty
    \item Chest pain, choking, or signs of stroke
    \item Severe burns or suspected spinal injury
    \item Any situation in which the person's life or limb is at risk
\end{itemize}

\section*{Diagnosis and Investigations}
\begin{itemize}
    \item \textbf{Kit audit:} regularly check contents against a list, replace used and expired items, and ensure medicines are in date
    \item \textbf{Location check:} the kit must be in an easily found, known place, and household members should know where it is
    \item \textbf{Training check:} knowing how to use the kit's contents is more important than the contents themselves
    \item \textbf{Specific supplies:} workplaces and remote areas may need additional items such as eye wash and snake-bite equipment
    \item \textbf{Medicines review:} check expiry dates and ensure prescribed items such as EpiPens are present
\end{itemize}

\section*{Management}
\begin{itemize}
    \item \textbf{Stock the basics:} adhesive bandages of various sizes, sterile gauze, antiseptic wipes, bandages including compression bandages, triangular bandages, tweezers, scissors, safety pins, gloves, a thermometer, paracetamol and ibuprofen, antihistamines, a cold pack, and a first aid guide
    \item \textbf{Add setting-specific items:} cling film for burns, an eye wash, a torch, a blanket, and snake-bite bandages for Australian settings
    \item \textbf{Store properly:} in a dry, accessible place, out of reach of small children but known to adults
    \item \textbf{Maintain the kit:} replace used items immediately, check expiry dates every few months, and refresh the kit after the first aid course guidelines change
    \item \textbf{Learn to use it:} complete a first aid course so the kit's contents are familiar
    \item \textbf{Personalise:} include medicines and equipment for family members' conditions
    \item \textbf{Workplace compliance:} follow the regulated first aid requirements for the workplace
\end{itemize}

\section*{Complications}
\begin{itemize}
    \item Treating serious injuries with a kit instead of calling for help
    \item Using expired medicines, dressings, or equipment
    \item Not being able to find items quickly in an emergency
    \item Infection from using unsterile or contaminated items
    \item Giving medicines without checking doses or allergies
    \item False confidence from having a kit without first aid training
\end{itemize}

\section*{Prognosis and Outcomes}
A well-stocked, maintained first aid kit, combined with basic training, improves outcomes across the full range of minor and moderate injuries, reducing infection, controlling bleeding, and managing pain until professional care. For serious emergencies, the kit's real value is in the first minutes, and its contents are secondary to calling 000 and performing CPR. Households that check their kit regularly and refresh their skills are prepared for the injuries most likely to occur, and their preparedness makes emergencies calmer and safer.

\section*{Prevention and Self-Care}
\begin{itemize}
    \item Assemble or buy a kit matched to your setting and check it monthly
    \item Keep it in one known place and tell family members and visitors
    \item Replace used and expired items immediately
    \item Complete and refresh a first aid course
    \item Keep emergency numbers visible, including 000 and 13 11 26 for poisons
    \item Add personal medicines and equipment for household conditions
    \item For workplaces, follow the mandatory first aid code of practice
    \item Store the kit out of children's reach but not locked away
\end{itemize}

\section*{Sources and Further Reading}
The following reputable sources provide further information for healthcare students and patients seeking reliable, current advice about first aid kits.
\begin{itemize}
    \item St John Ambulance Australia. \textit{First aid kits.} https://stjohn.org.au/
    \item Australian Red Cross. \textit{First aid kit contents.} https://www.redcross.org.au/
    \item Safe Work Australia. \textit{First aid in the workplace.} https://www.safeworkaustralia.gov.au/
    \item Healthdirect Australia. \textit{First aid kit checklist.} https://www.healthdirect.gov.au/
    \item NHS. \textit{Your first aid kit.} https://www.nhs.uk/
    \item World Health Organization. \textit{Emergency preparedness.} https://www.who.int/
\end{itemize}
